\chapter{La Torre: la campaña semana a semana}
\label{cap:s-torre}
\textit{Materias: Grandes volúmenes (datos que llegan) y Mercadotecnia (qué hacer si la afluencia rebasa la capacidad).
Pregunta: ¿qué hace la campaña esta semana?}

\section{La idea}
El plan de dinero se arma por mes, pero el mundo cambia cada semana: pasa una tormenta, llueve fuera de lo normal, llega
más gente de la esperada o el norte se llena. La Torre vigila esas señales cada semana y, si hace falta, \textbf{pausa} un
anuncio o \textbf{enciende} otro, con reglas fijadas de antemano.

\begin{intuicion}
Es el copiloto del viaje. El plan dice ``vamos por la carretera federal'', pero si el copiloto ve un choque adelante,
cambia de ruta sin tener que rehacer todo el viaje. Las reglas para cambiar de ruta se acordaron antes de salir.
\end{intuicion}

\section{Las cuatro señales}
\begin{center}
\small
\begin{tabularx}{\linewidth}{>{\raggedright\arraybackslash}p{0.25\linewidth}YY}
\encabezado \h{Señal} & \h{De dónde sale} & \h{Qué provoca} \\
\textbf{Tormenta cerca} & Trayectorias de la NOAA, con la regla de 200 km & Pausa el anuncio de ese lugar esa semana \\
\filagris\textbf{Clima raro} & Isolation Forest sobre la lluvia, el viento y el calor de la semana & Pausa \\
\textbf{Llegó más gente de la esperada} & El dato real del mes contra el rango del 90\,\% del Pronóstico & Pausa \\
\filagris\textbf{El norte se llenó} & Ocupación semanal de Cancún y la Riviera Maya (SECTUR) & Enciende ``¿Ibas al norte?'' \\
\bottomrule
\end{tabularx}
\normalsize
\end{center}

Si llega \textbf{menos} gente de la esperada, no se pausa nada: eso es espacio, justo lo que la campaña busca llenar.

\section{Las decisiones de Brandon}
\begin{itemize}
\item \textbf{Qué hacer en el sur con una mala semana}: ``pausar esa semana y guardar el dinero''. El dinero no se pierde:
  pasa a la siguiente semana permitida del mismo lugar. Se descartó bajar el anuncio a la mitad (se seguiría anunciando con
  mal clima) y solo avisar (no cuidaría nada).
\item \textbf{Qué hacer si el norte se llena}: encender el anuncio ``¿Ibas al norte?'', que invita al lugar del sur con más
  espacio ese mes. ``Lleno'' es el mismo corte que ya usa el Radar: 85.92\,\% de ocupación.
\item \textbf{Cómo mostrarlo}: como una reproducción grabada en la página (GitHub Pages no puede transmitir en vivo) y,
  en la computadora del proyecto, en vivo con el servidor.
\item \textbf{Qué es ``clima raro''}: más raro que el 95\,\% de las semanas que el modelo ya conoce.
\end{itemize}

\begin{error}
\textbf{El modelo de clima raro marcaba casi la mitad de las semanas.} El corte automático de Isolation Forest señalaba como
raras el \textbf{47\,\%} de las semanas de Chetumal, lo que habría pausado media campaña. Dos causas: con pocas semanas, el
modelo no distinguía bien, y desde 2022 llueve más y hace más calor que en 2019–2021, así que todo lo reciente le parecía
raro. Se cambió a un modelo por año que aprende de todos los años anteriores, y Brandon eligió el corte del 95\,\%. Quedaron
\textbf{21 semanas raras en Chetumal} y 17 en Kohunlich, entre 7\,\% y 9\,\% de las semanas.
\end{error}

\section{El dinero que no se pierde}
\begin{ejemplo}[ · una pausa]
Si un mes tiene 4 lunes y \$4,000 de presupuesto, cada semana le tocan \$1,000. Si la primera semana se pausa, esos
\$1,000 se guardan. La segunda semana, si todo está bien, se gastan \$2,000: los suyos y los guardados.
\end{ejemplo}

\section{Lo que pasó al reproducir 239 semanas reales}
Se revivieron \textbf{239 semanas reales}, del 3 de enero de 2022 al 27 de julio de 2026, como si llegaran una por una, con
Spark Structured Streaming (capítulo~\ref{cap:s-cajas}). El plan de dinero se aplicó a esos años para ver qué habría pasado.

\begin{center}
\small
\begin{tabular}{lrrrr}
\encabezado \h{Lugar} & \h{Encendido} & \h{Pausado} & \h{Temporada alta} & \h{Fuera del plan} \\
Chetumal & 145 & 19 & 18 & 57 \\
Bahía Calderitas–Oxtankah & 106 & 14 & 62 & 57 \\
Ruta de las pirámides & 104 & 15 & 63 & 57 \\
\bottomrule
\end{tabular}
\normalsize
\end{center}
(``Fuera del plan'' son julio, agosto y septiembre: el plan solo cubre de octubre a junio. En esas semanas la Torre vigila
pero no gasta.)

\begin{itemize}
\item \textbf{48 pausas} en total. Algunas tuvieron más de un motivo: 44 por clima raro, 9 por tormenta cerca (Lisa en 2022,
  Nadine y Sara en 2024) y 5 porque llegó más gente de la esperada a Chetumal en mayo de 2025.
\item \textbf{No se perdió dinero}: los \$876,901 planeados en las semanas con plan se gastaron todos.
\item \textbf{``¿Ibas al norte?'' se encendió 6 semanas}, todas hacia la Bahía. El norte casi nunca se satura en una semana
  suelta (Cancún 1 semana, la Riviera Maya 9): su señal fuerte es la temporada, no la semana.
\item \textbf{La prueba de independencia}: el modelo de clima raro marcó las tres semanas de tormenta sin saber que hubo
  tormenta.
\end{itemize}

\begin{teprofe}
\emph{``¿Esto es en tiempo real?''} Es una \textbf{reproducción de datos reales}: las semanas pasadas se procesan en orden,
una por una, con la misma tecnología que procesaría datos que llegan en vivo. Si mañana existiera una fuente semanal del
sur, se cambia la carpeta de entrada y las reglas siguen iguales. En la página se dice así: ``reproducción de datos
históricos reales''.
\end{teprofe}

\section{Lo que se declara}
\begin{itemize}
\item El sur no tiene ningún dato oficial semanal; la Torre del sur usa el clima, las tormentas y el dato mensual.
\item Las tormentas de 2026 todavía no se publican: esas semanas dicen ``sin dato'', no ``sin tormenta''.
\item Se supone que el dato del mes se conoce al empezar el mes siguiente; en la realidad tarda uno o dos meses.
\end{itemize}
